\section{Related Work}
\label{sec:related_work}

\subsection{Graph Representation Learning for Anti-Money Laundering \& Financial Fraud}
\label{sec:related_graph_aml}
The application of Graph Neural Networks (GNNs) to financial forensic analysis was pioneered by \citet{weber2019anti} on the Elliptic Bitcoin benchmark, demonstrating that spatial architectures such as Graph Convolutional Networks (GCN) \citep{kipf2017semi} and GraphSAGE \citep{hamilton2017inductive} capture localized topological connectivity patterns beyond isolated tabular classifiers. Subsequent research by \citet{bellei2024elliptic2} introduced the Elliptic2 multi-asset benchmark, highlighting higher-order relational subgraphs. In banking networks, \citet{cheng2023anti} introduced group-aware deep graph learning to capture collective community-level transaction behavior, while \citet{fu2025nowhere} developed H$^2$IDE to untangle fraudulent camouflage on multi-relational graphs via disentangled homophily and heterophily identification. More broadly, \citet{qiao2025deep} surveyed modern deep graph anomaly detection paradigms, highlighting the persistent challenges of structural camouflage, degree skew, and non-uniform relation topologies.

Graph representations have also been deployed for Ethereum blockchain forensics \citep{zheng2020xblock}, multi-bank synthetic wire simulations \citep{suzumura2019amlsim, lorenz2020samld}, streaming credit card fraud detection \citep{dalpozzolo2018credit}, and cross-institutional federated learning architectures \citep{wang2023automated}. However, prior AML graph frameworks largely operate over static topological snapshots. Consequently, they remain insensitive to continuous multi-scale temporal velocity patterns and vulnerable to severe feature over-smoothing when money launderers attach synthetic chaff edges to high-degree commercial hubs \citep{dou2020enhancing}.

\subsection{Continuous-Time Dynamic Temporal Graphs}
\label{sec:related_temporal_graphs}
Financial transaction networks evolve through asynchronous, continuous-time transaction events. Discrete snapshot dynamic GNNs, such as EvolveGCN \citep{pareja2020evolvegcn}, discretize event streams into regular time slices, discarding intra-slice temporal ordering and obscuring sub-second velocity bursts. While continuous frameworks including Temporal Graph Networks (TGN) \citep{rossi2020temporal}, TGAT \citep{xu2020inductive}, and continuous anomaly detectors \citep{wang2024temporal} preserve exact event order, standard single exponential decay kernels ($\exp(-\lambda \Delta t)$) rapidly attenuate to zero over multi-week dormant holding periods. In operational financial crime surveillance, single-decay dynamic GNNs struggle to simultaneously model sub-second structuring bursts alongside multi-week holding intervals characteristic of layering chains.

\subsection{Adversarial Graph Learning \& Topological Camouflage Defense}
\label{sec:related_adversarial}
Adversarial attacks on graph neural networks manipulate connectivity or node attributes to induce misclassification \citep{zugner2018adversarial}. In financial forensics, adversaries employ topological camouflage: laundering entities deliberately attach transactional chaff edges to high-degree, reputable commercial hubs (e.g., payment gateways, exchanges, utility merchants) to blend into background traffic \citep{dou2020enhancing}. Because standard graph convolutional aggregation treats incident edges with uniform or degree-normalized weights, camouflage connections drag illicit node embeddings toward majority benign centroids.

To counter camouflage, CARE-GNN \citep{dou2020enhancing} introduced reinforcement learning to select informative neighbors on multi-relational graphs, while \citet{ou2025camfd} developed CamFD to detect temporal behavioral inconsistencies of camouflaged fraudsters on dynamic graphs. However, existing camouflage defenses either rely on static relation heuristics or assume homophilic interactions, lacking continuous-time soft-gating that dynamically prunes high-degree banking chaff while providing provable finite-sample risk coverage.

\subsection{Imbalanced Graph Representation Learning \& Latent Oversampling}
\label{sec:related_imbalance}
Real-world financial transaction graphs exhibit extreme class imbalance, often with illicit prevalence $\pi < 0.05\%$. Traditional Euclidean oversampling techniques such as SMOTE \citep{chawla2002smote} operate strictly on tabular feature vectors, generating synthetic instances without topologically valid relational edges. To bridge topological oversampling, \citet{zhao2021graphsmote} proposed GraphSMOTE, performing feature interpolation in the intermediate latent embedding space and decoding synthetic edges via an auxiliary link generator. However, standard GraphSMOTE executes unconstrained latent interpolation across the global minority class pool, neglecting distinct structural laundering typologies (cyclic wash trading vs. smurfing fan-out vs. pass-through layering), risking synthetic instances that violate domain-level flow-balance regularity.

\subsection{Conformal Risk Control, Selective Classification \& Temporal Drift}
\label{sec:related_conformal}
Conformal prediction, formalized by \citet{vovk2005algorithmic} and extended by \citet{bates2021distribution} and \citet{angelopoulos2021gentle}, provides distribution-free, finite-sample uncertainty quantification. Recent research has investigated conformal prediction on GNNs \citep{huang2024conformal} and class-conditional conformal prediction under class imbalance \citep{ding2023class}. While standard marginal conformal prediction guarantees valid overall coverage ($1-\alpha$), it frequently suffers from severe conditional under-coverage on the rare class in highly imbalanced datasets. Under temporal non-stationarity, standard exchangeability is violated and coverage must be explicitly bounded through distribution drift tracking, such as Adaptive Conformal Inference (ACI) \citep{gibbs2021adaptive}. Unlike heuristic uncertainty frameworks, C-STGB leverages distribution-free Conformal Risk Control to deliver rigorous finite-sample class-conditional coverage bounds ($\ge 1 - \alpha$), backed by Martingale convergence under ACI tracking, while delineating the boundary where GNN message passing is indispensable vs. where tabular trees remain competitive.
